\subsection{Hardware Sampling Implementation}

I have to differentiate the hws internal implementation targeting \glspl{cpu} and \glspl{gpu}.
To sample information on the supported \glspl{gpu}, I use the respective vendor-provided library. 
For NVIDIA \glspl{gpu} this is the \gls{nvml} library, for AMD \glspl{gpu} it is the \gls{rocmsmi} library, and for Intel \glspl{gpu} it is the Level Zero library. 
While the \glspl{api} are relatively similar for \gls{nvml} and \gls{rocmsmi}, Level Zero has a significantly different approach, which can also be seen in \autoref{code:vendor_library_calls}. 
Nevertheless, all functions are simple library calls.

On the \gls{cpu}, however, this is different since there is no simple library that I can call to gather all the necessary information. 
Here, I use three different command line utilities: to retrieve general \gls{cpu} information, like core-count, socket-count, cache-sizes, etc., I use \texttt{lscpu}~\cite{qian_lscpu_nodate}, for memory-related information, i.e., available, used, and total main memory, I use \texttt{free}~\cite{small_free_nodate}, and for everything else like frequency or power related information I use \texttt{turbostat}~\cite{brown_turbostat_nodate}. 
In all three cases, I create a new subprocess, call the respective utility, and then return the standard output as a string. 
This string contains all necessary information and is then post-processed to extract the data inside the \gls{cpu} hardware sampler. 
For \texttt{lscpu}, the subprocess simply calls \texttt{lscpu}. 
In the case of the \texttt{free} command, the command line flag \texttt{-b} is added to be sure that the memory size information is output in bytes. 
For \texttt{turbostat} it is a bit different. 
First, I have to note that to get the full \texttt{turbostat} information, it has to be called with root privileges. 
However, in newer \texttt{turbostat} versions, it is possible to get a subset of all possible information even if \texttt{turbostat} is not called with root privileges. 
I handle that distinction directly in CMake: it tests whether \texttt{turbostat} can be executed with root privileges and if this is not the case it tests whether it can be executed without root privileges. 
In the latter case, gathering information via \texttt{turbostat} is completely disabled. 
In the former case, I sample the information using the \texttt{turbostat} command \texttt{sudo turbostat -n 1 -i 0.001 -S -q} with \texttt{sudo} only present if \texttt{turbostat} can be executed with root privileges. 
This command samples all available information immediately (\texttt{-i 0.001}) exactly once (\texttt{-n 1}) and outputs the information in a short summary format (\texttt{-S}) skipping potential header information (\texttt{-q}). 
An example \texttt{turbostat} output with the mentioned command looks like displayed in \autoref{code:turbostat_output}. 

\begin{listing}
    \begin{minted}[fontsize=\setfontsize{6.5}]{text}
Avg_MHz  Busy%  Bzy_MHz  TSC_MHz  IPC   IRQ  POLL  C1   C2   POLL%  C1%   C2%     CorWatt  PkgWatt
39       1.77   2212     4686     0.92  818  1     156  819  0.01   3.03  110.73  4.58     264.06
    \end{minted}
    \caption{%
    An example output of the command \texttt{sudo turbostat -n 1 -i 0.001 -S -q} with root privileges enabled for \texttt{turbostat} on an AMD EPYC 9274F \gls{cpu}.
    }
    \label{code:turbostat_output}
\end{listing}
Note that the output of \texttt{turbostat} also depends on the used \gls{cpu}. 
For example, the output \autoref{code:turbostat_output} was gathered on an AMD EPYC 9274F \gls{cpu} and consists of $14$ table entries. 
If I run the same command for example on an Intel i5-10210U \gls{cpu}, the output consists of $52$ columns mostly adding more idle state information but also temperature-related information missing for the AMD \gls{cpu}.
